\documentclass[12pt]{article}
\usepackage[utf8]{inputenc}
\usepackage[T1]{fontenc}
\usepackage[margin=1in]{geometry}
\usepackage{setspace}
\usepackage{xcolor}
\usepackage[colorlinks=true, urlcolor=blue!60!black, linkcolor=blue!60!black]{hyperref}
\setlength{\parindent}{2em}
\setlength{\parskip}{1.25em}
\renewcommand{\baselinestretch}{1.0}
\begin{document}
\begin{center}
{\LARGE \textit{West Virginia v. B. P. J. / Little v. Hecox}, Nos. 24-43 \& 24-38 (2026)\par}
\vspace{0.8em}
{\large Case Brief\par}
\end{center}
\vspace{1em}

\section{Memory Jogger}

In consolidated cases decided June 30, 2026, the Supreme Court held 9--0 in judgment under Title IX and 6--3 under the Equal Protection Clause that states may restrict public-school women's and girls' sports teams to biological females without making individualized exceptions for transgender girls and women (\href{https://www.courtlistener.com/opinion/10882186/west-virginia-v-b-p-j/}{CourtListener}; \href{https://www.law.cornell.edu/supremecourt/text/24-43}{Cornell LII}).

\section{Detailed Case Facts}

Between 2020 and 2026, 27 states enacted statutes restricting participation on public-school women's and girls' athletic teams to biological females, alongside parallel policies adopted in 2025--2026 by the NCAA, the U.S. Olympic and Paralympic Committee (USOPC), and the International Olympic Committee (IOC) (\href{https://www.law.cornell.edu/supremecourt/text/24-43}{Cornell LII}). These consolidated cases challenged two such laws from West Virginia and Idaho (\href{https://www.law.cornell.edu/supremecourt/text/24-43}{Cornell LII}).

In \textit{West Virginia v. B. P. J.} (No. 24-43), the West Virginia Legislature enacted the Save Women's Sports Act (H.B. 3293) in 2021, requiring public secondary schools and state colleges to designate athletic teams as male, female, or coed based on ``biological sex'' (defined by reproductive biology and genetics at birth) and prohibiting biological males from competing on female teams in contact sports or sports selected by competitive skill, W. Va. Code Ann. \S\S~18-2-25d(a)--(c) (\href{https://www.law.cornell.edu/supremecourt/text/24-43}{Cornell LII}; \href{https://harvardlawreview.org/print/vol-140/west-virginia-v-b-p-j-ex-rel-jackson/}{Harvard Law Review}). Respondent B. P. J. (Becky Pepper-Jackson) was assigned male at birth, socially transitioned as a girl in third grade, was diagnosed with gender dysphoria around age nine, began puberty-blocking medication to prevent endogenous male puberty, and started estrogen therapy in sixth grade (\href{https://www.law.cornell.edu/supremecourt/text/24-43}{Cornell LII}; \href{https://www.scotusblog.com/2026/06/court-rules-that-states-can-exclude-transgender-athletes-from-girls-and-womens-sports-teams/}{SCOTUSblog}). When B. P. J. prepared to enter middle school, her principal informed her mother, Heather Jackson, that the Act barred B. P. J. from joining the girls' cross-country and track-and-field teams (\href{https://www.law.cornell.edu/supremecourt/text/24-43}{Cornell LII}). During subsequent litigation under court-ordered injunctions, B. P. J. competed on girls' teams and later won the West Virginia Class AAA high school state championship in girls' shot put and placed fourth in girls' discus (\href{https://www.law.cornell.edu/supremecourt/text/24-43}{Cornell LII}).

In \textit{Little v. Hecox} (No. 24-38), the Idaho Legislature enacted the Fairness in Women's Sports Act (H.B. 500) in 2020, similarly defining sex by biology, barring biological males from female school teams from primary school through college, and providing a medical verification mechanism for sex disputes, Idaho Code Ann. \S\S~33-6202 to 33-6203 (\href{https://www.law.cornell.edu/supremecourt/text/24-43}{Cornell LII}; \href{https://harvardlawreview.org/print/vol-140/west-virginia-v-b-p-j-ex-rel-jackson/}{Harvard Law Review}). Respondent Lindsay Hecox, a transgender woman who underwent male puberty and began hormone therapy in college, attended Boise State University, competed on the women's club soccer team, and sought to try out for the women's Division I track and cross-country teams (\href{https://www.law.cornell.edu/supremecourt/text/24-43}{Cornell LII}).

\section{Procedural History}

In \textit{West Virginia v. B. P. J.}, B. P. J. (through her mother) sued state and county education officials in the U.S. District Court for the Southern District of West Virginia, asserting claims under Title IX of the Education Amendments of 1972, 20 U.S.C. \S~1681(a), and the Equal Protection Clause of the Fourteenth Amendment (\href{https://www.law.cornell.edu/supremecourt/text/24-43}{Cornell LII}). Judge Goodwin initially entered a preliminary injunction allowing B. P. J. to compete, 550 F. Supp. 3d 347 (S.D. W. Va. 2021), but later granted summary judgment to West Virginia on both claims and dissolved the injunction, 649 F. Supp. 3d 220 (S.D. W. Va. 2023) (\href{https://harvardlawreview.org/print/vol-140/west-virginia-v-b-p-j-ex-rel-jackson/}{Harvard Law Review}). The Fourth Circuit granted an injunction pending appeal (which the Supreme Court declined to vacate in 2023) and then reversed in a divided opinion, directing summary judgment for B. P. J. on Title IX and remanding the Equal Protection claim for further factfinding on whether B. P. J. possessed any athletic advantage, 98 F.4th 542 (4th Cir. 2024) (\href{https://www.law.cornell.edu/supremecourt/text/24-43}{Cornell LII}; \href{https://constitution.congress.gov/browse/essay/intro.9-4-7/ALDE_00000140/}{Constitution Annotated}).

In \textit{Little v. Hecox}, Hecox sued Governor Bradley Little and Idaho officials in the U.S. District Court for the District of Idaho, alleging an Equal Protection violation (\href{https://www.law.cornell.edu/supremecourt/text/24-43}{Cornell LII}). Chief Judge Nye preliminarily enjoined enforcement of H.B. 500, 479 F. Supp. 3d 930 (D. Idaho 2020), and the Ninth Circuit affirmed, 104 F.4th 1061 (9th Cir. 2024) (\href{https://harvardlawreview.org/print/vol-140/west-virginia-v-b-p-j-ex-rel-jackson/}{Harvard Law Review}). The Supreme Court granted certiorari in both cases in July 2025, 606 U.S. 959 (2025) (\href{https://www.law.cornell.edu/supremecourt/text/24-43}{Cornell LII}). In September 2025, Hecox filed a suggestion of mootness declaring she would cease playing covered women's sports at Boise State, which the Supreme Court rejected (\href{https://www.law.cornell.edu/supremecourt/text/24-43}{Cornell LII}). The Court heard consolidated oral argument on January 13, 2026, and issued its decision on June 30, 2026 (\href{https://www.law.cornell.edu/supremecourt/text/24-43}{Cornell LII}).

\section{Judicial Votes}

The Supreme Court reversed and remanded both appellate judgments, splitting \textbf{9--0 in judgment} (6--3 on rationale) on Title IX and \textbf{6--3} on Equal Protection and mootness (\href{https://www.law.cornell.edu/supremecourt/text/24-43}{Cornell LII}; \href{https://www.scotusblog.com/2026/06/court-rules-that-states-can-exclude-transgender-athletes-from-girls-and-womens-sports-teams/}{SCOTUSblog}):
\begin{itemize}
\item \textbf{Majority (6):} Justice Brett Kavanaugh (author), joined in full by Chief Justice John Roberts and Associate Justices Clarence Thomas, Samuel Alito, Neil Gorsuch, and Amy Coney Barrett.
\item \textbf{Concurring Opinions:} Justice Clarence Thomas (joined majority in full; filed solo concurrence) and Justice Neil Gorsuch (joined majority in full; filed solo concurrence).
\item \textbf{Concurring in the Judgment in Part and Dissenting in Part (3):} Justice Sonia Sotomayor (author), joined in full by Associate Justices Elena Kagan and Ketanji Brown Jackson (concurring in the judgment on B. P. J.'s Title IX claim; dissenting on B. P. J.'s Equal Protection claim and arguing \textit{Hecox} was moot).
\item \textbf{Separate Concurring in the Judgment in Part and Dissenting in Part (1):} Justice Ketanji Brown Jackson (joined Sotomayor's opinion in full; wrote separately to challenge the majority's narrow reading of ``sex'' in Title IX).
\end{itemize}

\section{Holding}

\textbf{Judgment:} Reversed and remanded in both No. 24-43 (Fourth Circuit) and No. 24-38 (Ninth Circuit) (\href{https://www.law.cornell.edu/supremecourt/text/24-43}{Cornell LII}).

\textbf{Rule Announced:} (1) Under Title IX (20 U.S.C. \S~1681(a)), the 1974 Javits Amendment (\S~844, 88 Stat. 612), and the longstanding 1975 implementing regulations (34 C.F.R. \S~106.41(b)), ``sex'' refers to biological sex, and schools may permissibly maintain separate women's and girls' sports teams limited to biological females (\href{https://www.law.cornell.edu/supremecourt/text/24-43}{Cornell LII}). (2) Under the Equal Protection Clause of the Fourteenth Amendment, state laws restricting women's and girls' sports teams to biological females classify on the basis of biological sex rather than transgender status and satisfy intermediate scrutiny because they are substantially related to the important governmental objectives of participant safety and competitive fairness; states are not constitutionally required to grant individualized or subclass-based exemptions for transgender athletes who take puberty blockers or cross-sex hormones (\href{https://www.law.cornell.edu/supremecourt/text/24-43}{Cornell LII}; \href{https://constitution.congress.gov/browse/essay/intro.9-4-7/ALDE_00000140/}{Constitution Annotated}).

\textbf{Holding vs. Dicta:} The \emph{holding} establishes that neither Title IX nor the Equal Protection Clause prohibits states from maintaining female school sports teams defined solely by biological sex, and that \textit{Hecox} was not moot because Hecox remained enrolled at Boise State after a post-certiorari shift in position (\href{https://www.law.cornell.edu/supremecourt/text/24-43}{Cornell LII}). By contrast, the Court expressly reserved---and treated as \emph{non-decided dicta/caveats}---(a) whether Title IX or Equal Protection \emph{permits} states or schools to allow transgender girls and women to compete on female teams, (b) participation by biological females on male or coed teams, (c) whether transgender status or gender identity triggers heightened scrutiny under the Equal Protection Clause (citing \textit{United States v. Skrmetti}, 605 U.S. 495 (2025)), and (d) how \textit{Bostock v. Clayton County}, 590 U.S. 644 (2020), applies to Title IX outside competitive athletics (\href{https://www.law.cornell.edu/supremecourt/text/24-43}{Cornell LII}).

\section{Analysis of Opinions}

\subsection{Majority Opinion (Justice Kavanaugh)}

Justice Kavanaugh opened by celebrating Title IX's transformative expansion of female athletics since 1972 and emphasizing that ``enduring'' physical differences between biological males and females in height, weight, strength, speed, endurance, and jumping ability necessitate sex-separated teams to ensure safety and competitive fairness (\href{https://www.law.cornell.edu/supremecourt/text/24-43}{Cornell LII}).

On \textbf{Title IX} (raised only by B. P. J.), Kavanaugh held that the ordinary meaning of ``sex'' in the 1972 statute, the 1974 Javits Amendment, and the contemporaneous 1975 HEW regulations---which carry special interpretive weight under \textit{Loper Bright Enterprises v. Raimondo}, 603 U.S. 369 (2024)---is biological sex (\href{https://www.law.cornell.edu/supremecourt/text/24-43}{Cornell LII}). Because B. P. J. conceded that Title IX allows separate teams excluding most biological males, the text provides no basis for mandating an exception for biological males who identify as female and take puberty blockers or hormones (\href{https://www.law.cornell.edu/supremecourt/text/24-43}{Cornell LII}). Drawing a bright biological line is ``reasonable'' under the Javits Amendment because sports are ``generally zero sum'': every roster spot, starting role, or medal won by a biological male displaces a female athlete (\href{https://www.law.cornell.edu/supremecourt/text/24-43}{Cornell LII}). Kavanaugh rejected B. P. J.'s reliance on Title VII and \textit{Bostock}, noting that Title VII forbids workplace sex distinctions whereas Title IX expressly authorizes sex-separated sports teams (\href{https://www.law.cornell.edu/supremecourt/text/24-43}{Cornell LII}).

On \textbf{Equal Protection}, Kavanaugh first confirmed in footnote 4 that \textit{Hecox} was not moot under voluntary-cessation doctrine because Hecox remained enrolled at Boise State and post-certiorari maneuvers to insulate a lower-court judgment merit skepticism (\href{https://www.law.cornell.edu/supremecourt/text/24-43}{Cornell LII}). On the merits, he treated both statutes as sex-based classifications triggering intermediate scrutiny and held them substantially related to the important state interests in safety and competitive fairness (\href{https://www.law.cornell.edu/supremecourt/text/24-43}{Cornell LII}). Relying on \textit{Tuan Anh Nguyen v. INS}, 533 U.S. 53 (2001), \textit{Michael M. v. Superior Court}, 450 U.S. 464 (1981), and First Amendment intermediate-scrutiny precedents (\textit{Ward v. Rock Against Racism} and \textit{Edge Broadcasting}), Kavanaugh rejected the plaintiffs' as-applied challenge as a ``backdoor argument for strict scrutiny'' that would plunge courts into an unadministrable, intrusive quagmire of athlete-by-athlete physiological comparisons (\href{https://www.law.cornell.edu/supremecourt/text/24-43}{Cornell LII}). Relying on \textit{Skrmetti}, he held that the laws classify by biological sex rather than transgender status, would survive either rational-basis or intermediate scrutiny regardless, and warrant legislative deference amid ``ongoing medical and scientific debate'' (\href{https://www.law.cornell.edu/supremecourt/text/24-43}{Cornell LII}).

\subsection{Concurring Opinion (Justice Thomas)}

Justice Thomas joined the majority opinion in full and wrote separately to make two categorical points. First, he argued that transgender status is not a suspect or quasi-suspect class warranting heightened scrutiny because gender dysphoria is ``a mutable mental state that is the object of psychiatric treatment,'' akin to traits reviewed solely for a rational basis (\href{https://www.law.cornell.edu/supremecourt/text/24-43}{Cornell LII}). Second, he insisted that sex is immutable and binary, asserting that ``[m]en and boys with gender dysphoria are not women or girls, even if they believe that they are,'' and quoting philosopher Josef Pieper to argue that using language to obscure biological reality is to lie to the public (\href{https://www.law.cornell.edu/supremecourt/text/24-43}{Cornell LII}).

\subsection{Concurring Opinion (Justice Gorsuch)}

Justice Gorsuch---the author of \textit{Bostock}---joined the majority opinion in full and added two statutory points on Title IX. First, because Title IX was enacted under the Spending Clause, Congress must ``clearly and unambiguously'' give funding recipients notice of conditions attached to federal funds under \textit{Pennhurst State School and Hospital v. Halderman}, 451 U.S. 1 (1981), and \textit{Medina v. Planned Parenthood South Atlantic}, 606 U.S. 357 (2025); nothing in Title IX clearly alerts schools that single-sex biological sports teams are forbidden (\href{https://www.law.cornell.edu/supremecourt/text/24-43}{Cornell LII}). Second, Gorsuch explained that \textit{Bostock} is fully consistent with the Court's ruling: both cases read ``sex'' as biological sex and treat but-for causation similarly, but whereas firing an employee because of sex is indisputably unlawful ``discrimination'' under Title VII, maintaining separate-sex sports teams is expressly authorized---and not unlawful discrimination---under Title IX (\href{https://www.law.cornell.edu/supremecourt/text/24-43}{Cornell LII}).

\subsection{Opinion Concurring in the Judgment in Part and Dissenting in Part (Justice Sotomayor)}

Justice Sotomayor, joined by Justices Kagan and Jackson, agreed in Part III that B. P. J.'s Title IX claim failed solely because B. P. J. conceded that ``sex'' means biological sex in the sports regulations, while stressing that \textit{Bostock}'s causation framework still protects transgender students under Title IX in non-athletic school settings such as classes, clubs, admissions, and harassment (\href{https://www.law.cornell.edu/supremecourt/text/24-43}{Cornell LII}). She also contended that \textit{Hecox} should have been dismissed as moot under \textit{Munsingwear} because Hecox was nearing graduation and swore never to play covered sports again (\href{https://www.law.cornell.edu/supremecourt/text/24-43}{Cornell LII}).

On B. P. J.'s \textbf{Equal Protection} claim, Sotomayor dissented sharply, arguing that the Court should have affirmed the Fourth Circuit's remand for trial-court factfinding (\href{https://www.law.cornell.edu/supremecourt/text/24-43}{Cornell LII}). She faulted the majority for reviving rejected dissents from \textit{Caban v. Mohammed}, 441 U.S. 380 (1979), and \textit{United States v. Virginia}, 518 U.S. 515 (1996), and for ignoring precedents (\textit{Caban}, \textit{Virginia}, \textit{Trimble v. Gordon}, \textit{Sessions v. Morales-Santana}) requiring courts to examine whether a categorical sex ban unnecessarily excludes a discrete, readily identifiable subclass---transgender girls who never underwent endogenous male puberty and receive hormone therapy---for whom the state's interests may not be furthered (\href{https://www.law.cornell.edu/supremecourt/text/24-43}{Cornell LII}). Sotomayor exposed what she called the majority's ``contorted logic'': crediting the legislature's fear that transgender girls pose an existential threat to girls' sports while declaring the same subclass too small to merit equal protection scrutiny (\href{https://www.law.cornell.edu/supremecourt/text/24-43}{Cornell LII}). She further accused the majority of collapsing intermediate scrutiny into rational-basis deference by relying on rational-basis cases (\textit{Murgia}, \textit{Skrmetti}), citing extra-record post-argument studies, and overlooking that the NCAA and USOPC changed policies after Executive Order 14201 threatened federal funding (\href{https://www.law.cornell.edu/supremecourt/text/24-43}{Cornell LII}).

\subsection{Opinion Concurring in the Judgment in Part and Dissenting in Part (Justice Jackson)}

Justice Jackson joined Justice Sotomayor's opinion in full and wrote separately to criticize the majority for unnecessarily holding that ``sex'' in Title IX means exclusively biological sex (\href{https://www.law.cornell.edu/supremecourt/text/24-43}{Cornell LII}). While B. P. J.'s concession resolved her Title IX claim here, Jackson argued that the Court's sex-stereotyping precedents (\textit{Price Waterhouse v. Hopkins}, \textit{Manhart}, \textit{Davis}, and \textit{Jackson v. Birmingham Board of Education}) demonstrate that Title IX protects individuals' gender presentation and gender identity against institutional gender-based expectations, not merely sex assigned at birth (\href{https://www.law.cornell.edu/supremecourt/text/24-43}{Cornell LII}).

\section{Examples: Future Applications}

\begin{enumerate}
\item \textbf{Same-Side (Upheld --- State Ban on Transgender Women in Public University Intramural Basketball):} A state enacts a law restricting all public university women's intramural and club basketball leagues to biological females, and a transgender female student taking estrogen sues under Title IX and the Equal Protection Clause. \emph{Why:} Under \textit{B. P. J.}, basketball is a contact sport selected by competitive skill, and states may draw a categorical biological-sex line for women's school athletics without conducting individualized physiological assessments.

\item \textbf{Same-Side (Upheld --- High School Athletic Association Rule for Girls' Swimming):} A state high school athletic association adopts a biological-sex eligibility rule for girls' varsity swimming, and a transgender girl who took puberty blockers challenges the rule as applied to her. \emph{Why:} \textit{B. P. J.} directly forecloses as-applied Equal Protection challenges seeking individualized exemptions from biological-sex sports classifications and confirms that Title IX's 1975 regulations permit such lines.

\item \textbf{Opposite-Side (Invalidated --- Public School Exclusion of Transgender Student from Coed Chess Club or Debate Team):} A school district bars a transgender student from joining its coed varsity chess club and debate team because the student is transgender, citing \textit{B. P. J.} \emph{Why:} Both Kavanaugh's majority opinion (expressly distinguishing a coed team that bars transgender students) and Sotomayor's and Gorsuch's opinions confirm that excluding transgender students from coed or non-athletic educational activities classifies based on transgender status/sex outside Title IX's athletic carve-out and violates Title IX under \textit{Bostock} and Equal Protection principles (\href{https://www.law.cornell.edu/supremecourt/text/24-43}{Cornell LII}).

\item \textbf{Opposite-Side (Invalidated --- Public School Expulsion or Harassment of a Transgender Student):} A federally funded public high school expels a student solely for socially transitioning and living as a girl, arguing that \textit{B. P. J.} holds Title IX does not protect transgender students. \emph{Why:} As Justice Gorsuch's concurrence and Justice Sotomayor's opinion make clear, \textit{B. P. J.} turned on Title IX's specific authorization of sex-separated sports teams under the Javits Amendment; expelling a student for being transgender is unlawful discrimination ``on the basis of sex'' under \textit{Bostock}'s causation logic (\href{https://www.law.cornell.edu/supremecourt/text/24-43}{Cornell LII}).

\item \textbf{Fence-Sitter (Borderline --- State Law Requiring Schools to Allow Transgender Girls on Girls' Teams Challenged by Cisgender Athletes):} A state passes an inclusive statute guaranteeing transgender girls who suppress testosterone the right to compete on public-school girls' sports teams, and cisgender female athletes sue under Title IX and the Equal Protection Clause arguing the inclusive policy deprives them of ``equal athletic opportunity.'' \emph{Why:} Footnote 2 of Kavanaugh's majority opinion explicitly notes that \textit{B. P. J.} does \emph{not} decide whether Title IX or the Equal Protection Clause permits (or forbids) schools to include transgender girls on female teams, leaving that mirror-image litigation open in the lower courts (\href{https://www.law.cornell.edu/supremecourt/text/24-43}{Cornell LII}).
\end{enumerate}

\section{Critique}

\textbf{Originalist and Textualist Perspectives:} Conservative, originalist, and textualist commentators---highlighted in Federalist Society analyses---praise \textit{B. P. J.} for adhering to the 1972 public meaning of ``sex'' in Title IX, enforcing Spending Clause clear-notice limits under \textit{Pennhurst}, and respecting democratic federalism on contested scientific and policy questions (\href{https://fedsoc.org/case/west-virginia-v-bpj}{Federalist Society}; \href{https://www.law.cornell.edu/supremecourt/text/24-43}{Cornell LII}). From this viewpoint, refusing to mandate judicially supervised hormone-and-physiology exemptions preserves workable bright-line rules and prevents courts from acting as national athletic commissions (\href{https://www.law.cornell.edu/supremecourt/text/24-43}{Cornell LII}).

\textbf{Progressive and Doctrinal Perspectives:} Progressive constitutional scholars raise three major critiques. First, Dean Erwin Chemerinsky warns that by omitting \textit{United States v. Virginia}'s ``exceedingly persuasive justification'' standard, rejecting individualized fit inquiries (\textit{Orr v. Orr}, \textit{Caban}), and importing rational-basis deference for scientific uncertainty, the majority dangerously watered down intermediate scrutiny for \emph{all} sex-discrimination claims (\href{https://www.scotusblog.com/2026/08/weakening-the-protection-against-sex-discrimination/}{SCOTUSblog --- Chemerinsky}; \href{https://www.abajournal.com/columns/article/chemerinsky-how-supreme-courts-decisions-had-an-impact-on-democracy}{ABA Journal}). Second, Professor Craig Konnoth argues that the litigation focused on the wrong issue: by conceding that ``sex'' meant biological sex and seeking an exception for a subclass of ``biological males,'' plaintiffs invited an unworkable comparator analysis rather than challenging whether states' statutory definitions of ``female'' (such as chromosomes or endogenous testosterone) actually track athletic performance (\href{https://www.scotusblog.com/2026/07/why-the-supreme-court-decided-the-wrong-issue-in-the-transgender-athletes-cases/}{SCOTUSblog --- Konnoth}). Third, the \textit{Harvard Law Review} Leading Case comment criticizes all five opinions for ignoring the coordinated nationwide political campaign behind anti-transgender bills and the resulting gap between Equal Protection animus doctrine and legislative deference (\href{https://harvardlawreview.org/print/vol-140/west-virginia-v-b-p-j-ex-rel-jackson/}{Harvard Law Review}).

\textbf{Underlying Competing Values:} The decision reflects a profound collision between competing conceptions of equality: safeguarding sex-separated competitive spaces and opportunities for biological females versus guaranteeing transgender youth individualized dignity, inclusion, and freedom from overbroad categorical exclusions, all set against the institutional question of whether legislatures or courts should resolve evolving scientific and cultural disputes.

\section{Key Quotations}

\begin{quote}
``Consistent with Title IX and the Equal Protection Clause, we hold that the States may maintain women's and girls' sports for biological females. They may determine eligibility for women's and girls' sports based on biological sex. The Constitution and Title IX do not require an overhaul of women's and girls' sports throughout America.''
\end{quote}
--- Justice Brett Kavanaugh, Majority Opinion, slip op. at 29 (\href{https://www.law.cornell.edu/supremecourt/text/24-43}{Cornell LII}; \href{https://www.courtlistener.com/opinion/10882186/west-virginia-v-b-p-j/}{CourtListener}).

\begin{quote}
``The underlying medical and scientific premise of the plaintiffs' entire equal protection argument is that at least some biological males who identify as female and take puberty blockers or hormones do not retain physical advantages over biological females. Even if true, that empirical claim would not alter the equal protection conclusion set forth above.''
\end{quote}
--- Justice Brett Kavanaugh, Majority Opinion, slip op. at 24 (\href{https://www.law.cornell.edu/supremecourt/text/24-43}{Cornell LII}; \href{https://www.courtlistener.com/opinion/10882186/west-virginia-v-b-p-j/}{CourtListener}).

\begin{quote}
``Nothing in Title IX clearly and unambiguously alerts funding recipients that they are prohibited from restricting a school-sponsored sports team to biological women or girls.''
\end{quote}
--- Justice Neil Gorsuch, Concurring Opinion (\href{https://www.law.cornell.edu/supremecourt/text/24-43}{Cornell LII}; \href{https://www.courtlistener.com/opinion/10882186/west-virginia-v-b-p-j/}{CourtListener}).

\begin{quote}
``Because of the Court's decision today, West Virginia, and any other state actor, can deny B. P. J. and others like her these experiences simply because it thinks they have an inherent athletic advantage, even if the facts show that they do not. In the end, to the Court, the facts do not matter, even though the consequences are serious.''
\end{quote}
--- Justice Sonia Sotomayor, Concurring in the Judgment in Part and Dissenting in Part (\href{https://www.law.cornell.edu/supremecourt/text/24-43}{Cornell LII}; \href{https://www.courtlistener.com/opinion/10882186/west-virginia-v-b-p-j/}{CourtListener}).

\begin{quote}
``In short, the majority is wrong to suggest that the term `sex' in Title IX `cannot plausibly be interpreted to refer to anything other than biological sex.' Title IX makes room for individuals to live in the gender they choose; it cares not just about sex assigned at birth but also about individuals' ability to match (or not) their gender presentation to their gender identity.''
\end{quote}
--- Justice Ketanji Brown Jackson, Concurring in the Judgment in Part and Dissenting in Part (\href{https://www.law.cornell.edu/supremecourt/text/24-43}{Cornell LII}; \href{https://www.courtlistener.com/opinion/10882186/west-virginia-v-b-p-j/}{CourtListener}).

\end{document}
